% ╔══════════════════════════════════════════════════════════════════════════╗
% ║  probe.tex --- PLANTILLA MAESTRA DE REPORTE                      ║
% ║  reporte_base.tex  |  Versión 1.0                                       ║
% ║  Motor: Tectonic   |  Preprocesador: Jinja2                             ║
% ╚══════════════════════════════════════════════════════════════════════════╝

\documentclass[11pt, a4paper, final]{article}


% ════════════════════════════════════════════════════════════════════════════
%  0. PAQUETES --- ORDENADOS POR DEPENDENCIA
% ════════════════════════════════════════════════════════════════════════════

% --- Codificación y tipografía ---
\usepackage[T1]{fontenc}
\usepackage[utf8]{inputenc}

% Fuente principal: Helvetica (o clon similar), ultra limpia y corporativa/industrial.
\usepackage{helvet}
\renewcommand{\familydefault}{\sfdefault}

% Fuente secundaria (para \texttt{}): IBM Plex Mono o Inconsolata para los datos crudos.
\usepackage[scale=0.95]{plex-mono} 
\usepackage{microtype}

% --- Idioma (es-noquoting: previene secuestro de < > por babel) ---
\usepackage[spanish, es-tabla, es-noquoting]{babel}

% --- Layout ---
\usepackage[
  a4paper,
  top=2.4cm, bottom=2.6cm,
  left=2.0cm, right=2.0cm,
  headheight=1.6cm,
  headsep=0.6cm,
  footskip=1.0cm
]{geometry}

% --- Color (cargado antes que los demás para que los usen) ---
\usepackage[table, dvipsnames, svgnames]{xcolor}

% --- Gráficos y diagramas ---
\usepackage{tikz}
\usepackage{pgfplots}
\pgfplotsset{compat=1.18}
\usetikzlibrary{
  calc,
  shapes.geometric,
  arrows.meta,
  positioning,
  decorations.pathreplacing,
  backgrounds,
  patterns,
  fit,
  shadows,
  babel
}

% --- Cajas de alerta y contenido ---
\usepackage{tcolorbox}
\tcbuselibrary{skins, breakable, theorems, hooks}

% --- Tablas avanzadas ---
\usepackage{booktabs}
\usepackage{array}
\usepackage{multirow}
\usepackage{tabularx}
\usepackage{colortbl}

% --- MAGIA PARA CENTRAR TABLAS ---
\renewcommand{\tabularxcolumn}[1]{>{\raggedright\arraybackslash}m{#1}}
\renewcommand{\arraystretch}{1.3}

% --- Encabezados y pies ---
\usepackage{fancyhdr}
\usepackage{lastpage}

% --- Matemáticas ---
\usepackage{amsmath}
\usepackage{amssymb}

% --- Utilidades de texto ---
\usepackage{parskip}
\usepackage{enumitem}
\usepackage{graphicx}
\usepackage{calc}
\usepackage{setspace}

% --- Hipervínculos (siempre al final) ---
\usepackage[
  hidelinks,
  pdfauthor={probe.tex},
  pdftitle={Informe de Diagnóstico Forense de Hardware},
  pdfsubject={Hardware Diagnostic Report}
]{hyperref}


% ════════════════════════════════════════════════════════════════════════════
%  1. PALETA DE COLORES CORPORATIVOS
% ════════════════════════════════════════════════════════════════════════════

% --- Escala primaria (azul marino / pizarra) ---
\definecolor{AMPrimary}   {HTML}{000000}
\definecolor{AMSecondary} {HTML}{333333}
\definecolor{AMAccent}    {HTML}{EE0000}
\definecolor{AMSlate}     {HTML}{555555}
\definecolor{AMLight}     {HTML}{F0F0F0}
\definecolor{AMBorder}    {HTML}{DDDDDD}
\definecolor{AMText}      {HTML}{111111}
\definecolor{AMCanvas}    {HTML}{FFFFFF}

% --- Colores de estado (semáforo diagnóstico) ---
\definecolor{StatusOK}     {HTML}{1B4332} % Verde táctico oscuro
\definecolor{StatusOKBg}   {HTML}{F0FFF4} % Blanco con tinte verde muy sutil
\definecolor{StatusOKMid}  {HTML}{2D6A4F} % Verde industrial para gráficos
\definecolor{StatusWarn}   {HTML}{FF5000}
\definecolor{StatusWarnBg} {HTML}{FFF4E6}
\definecolor{StatusWarnMid}{HTML}{FF8C00}
\definecolor{StatusCrit}   {HTML}{EE0000}
\definecolor{StatusCritBg} {HTML}{FDEEED}
\definecolor{StatusCritMid}{HTML}{FF3333}
\definecolor{StatusInfo}   {HTML}{888888}
\definecolor{StatusInfoBg} {HTML}{F5F5F5}
\definecolor{redtex}       {HTML}{FF0000} % Rojo vibrante terminal


% ════════════════════════════════════════════════════════════════════════════
%  2. CONFIGURACIÓN TIPOGRÁFICA Y SECCIONES
% ════════════════════════════════════════════════════════════════════════════

\usepackage[explicit]{titlesec}

% Sección: línea izquierda coloreada + numeración en azul
\titleformat{\section}
  {\large\bfseries\color{AMPrimary}}
  {}
  {0pt}
  {\textcolor{AMPrimary}{\rule[-0.2ex]{4pt}{1.2em}}\hspace{8pt}#1}
\titlespacing*{\section}{0pt}{14pt}{6pt}

% Subsección: línea base punteada
\titleformat{\subsection}
  {\normalsize\bfseries\color{AMSecondary}}
  {\thesubsection.}
  {6pt}
  {#1}
\titlespacing*{\subsection}{0pt}{10pt}{4pt}


% ════════════════════════════════════════════════════════════════════════════
%  3. CAJAS tcolorbox PERSONALIZADAS
% ════════════════════════════════════════════════════════════════════════════
%  Uso en el documento:
%    \begin{amboxok}{Título}   ... \end{amboxok}
%    \begin{amboxwarn}{Título} ... \end{amboxwarn}
%    \begin{amboxcrit}{Título} ... \end{amboxcrit}
%    \begin{amboxinfo}{Título} ... \end{amboxinfo}

% Macro interna para todas las cajas de estado
\newcommand{\amstateboxstyle}[5]{%
  % #1=colback  #2=colframe  #3=symbolo  #4=titulo  #5=contenido
  \begin{tcolorbox}[
    enhanced,
    colback=#1,
    colframe=#2,
    coltitle=#2,
    fonttitle=\bfseries\small,
    title={#3\enspace #4},
    boxrule=0.75pt,
    arc=0pt,
    left=8pt, right=8pt, top=5pt, bottom=5pt
  ]
    #5
  \end{tcolorbox}%
}

% PASS / OK --- Verde -> OPTIMAL
\newtcolorbox{amboxok}[1]{
  enhanced,
  colback=StatusOKBg, colframe=StatusOK,
  coltitle=white, fonttitle=\bfseries\small, % <--- COLTITLE BLANCO
  title={[ OPTIMAL ]\enspace #1},
  boxrule=0.75pt, arc=0pt,
  left=8pt, right=8pt, top=5pt, bottom=5pt
}

% WARNING --- Naranja -> DEGRADED
\newtcolorbox{amboxwarn}[1]{
  enhanced,
  colback=StatusWarnBg, colframe=StatusWarn,
  coltitle=white, fonttitle=\bfseries\small, % <--- COLTITLE BLANCO
  title={[ DEGRADED ]\enspace #1},
  boxrule=0.75pt, arc=0pt,
  left=8pt, right=8pt, top=5pt, bottom=5pt
}

% CRITICAL --- Rojo -> CRITICAL
\newtcolorbox{amboxcrit}[1]{
  enhanced,
  colback=StatusCritBg, colframe=StatusCrit,
  coltitle=white, fonttitle=\bfseries\small, % <--- COLTITLE BLANCO
  title={[ CRITICAL ]\enspace #1},
  boxrule=0.75pt, arc=0pt,
  left=8pt, right=8pt, top=5pt, bottom=5pt
}

% INFO --- Azul neutro -> UNKNOWN
\newtcolorbox{amboxinfo}[1]{
  enhanced,
  colback=StatusInfoBg, colframe=StatusInfo,
  coltitle=white, fonttitle=\bfseries\small, % <--- COLTITLE BLANCO
  title={[ UNKNOWN ]\enspace #1},
  boxrule=0.75pt, arc=0pt,
  left=8pt, right=8pt, top=5pt, bottom=5pt
}

% Caja de info de componente (banda izquierda, sin título flotante)
\newtcolorbox{componentinfo}{
  enhanced,
  colback=AMLight, colframe=AMPrimary,
  boxrule=0pt, leftrule=4pt,
  arc=0pt,
  left=10pt, right=8pt, top=5pt, bottom=5pt
}


% ════════════════════════════════════════════════════════════════════════════
%  4. CONFIGURACIÓN GLOBAL DE pgfplots
% ════════════════════════════════════════════════════════════════════════════

\pgfplotsset{
  % Estilo base aplicable con: \begin{axis}[probe.tex, ...]
  probe.tex/.style={
    % Ejes
    axis line style   = {color=AMSlate, thin},
    tick style        = {color=AMSlate, thin},
    tick label style  = {font=\scriptsize, color=AMText},
    label style       = {font=\footnotesize\bfseries, color=AMSlate},
    title style       = {font=\small\bfseries, color=AMPrimary, yshift=2pt},
    % Grilla
    grid              = both,
    grid style        = {line width=0.25pt, draw=AMBorder},
    major grid style  = {line width=0.4pt,  draw=AMBorder},
    minor grid style  = {line width=0.2pt,  draw=AMBorder!50},
    % Leyenda
    legend style = {
      font         = \scriptsize,
      draw         = AMBorder,
      fill         = AMCanvas,
      rounded corners = 0pt,
      inner sep    = 4pt,
      row sep      = 2pt,
    },
    % Fondo del área de plot
    axis background/.style = {fill=none},
    % Márgenes internos
    enlargelimits = 0.05,
  },
}


% ════════════════════════════════════════════════════════════════════════════
%  5. COMANDOS UTILITARIOS
% ════════════════════════════════════════════════════════════════════════════

% --- Badges de estado inline ---
% Uso: \badgeok  \badgefail  \badgewarn  \badgeinfo
\newcommand{\makebadge}[2]{%
  \tikz[baseline=-0.65ex, inner sep=0pt]\node[ %
    fill=#1, text=white,
    font=\bfseries\scriptsize,
    inner xsep=4pt, inner ysep=2pt,
    rounded corners=0pt
  ]{#2};%
}
\newcommand{\badgeok}  {\makebadge{StatusOK}{OPTIMAL}}
\newcommand{\badgefail}{\makebadge{StatusCrit}{CRITICAL}}
\newcommand{\badgewarn}{\makebadge{StatusWarn}{DEGRADED}}
\newcommand{\badgeinfo}{\makebadge{StatusInfo}{UNKNOWN}}

% --- Valor métrico destacado ---
% Uso: \mv{95$^\circ$C}  ->  valor en azul primario negrita
\newcommand{\mv}[1]{\textcolor{AMPrimary}{\textbf{#1}}}

% --- Separador horizontal de sección ---
\newcommand{\sectionrule}{%
  \vspace{4pt}%
  \noindent\textcolor{AMBorder}{\rule{\linewidth}{0.5pt}}%
  \vspace{4pt}%
}

% --- Fila de tabla coloreada alternada ---
\newcommand{\rowalt}{\rowcolor{AMLight}}

% --- Celda de encabezado de tabla ---
\newcommand{\thcell}[1]{\color{AMPrimary}\textbf{#1}}


% ════════════════════════════════════════════════════════════════════════════
%  6. ENCABEZADOS Y PIE DE PÁGINA (fancyhdr)
% ════════════════════════════════════════════════════════════════════════════

\pagestyle{fancy}
\fancyhf{}
\renewcommand{\headrulewidth}{0pt}
\renewcommand{\footrulewidth}{0pt}

% Encabezado: banda coloreada con regla inferior
\fancyhead[L]{%
  \colorbox{AMPrimary}{%
    \hspace{2pt}%
    \color{white}\small\textbf{probe\textcolor{redtex}{.tex}}%
    \hspace{4pt}%
  }%
  \textcolor{AMPrimary}{\rule[-4pt]{1pt}{14pt}}%
  \hspace{4pt}%
  \color{AMSlate}\small\textit{Diagnóstico Forense de Hardware}%
}
\fancyhead[R]{%
  \small\color{AMSlate}%
  << cliente_nombre >> \textcolor{AMBorder}{\textbullet}\ << fecha_reporte >>%
}
% Línea bajo encabezado
\renewcommand{\headrule}{%
  \vspace{-4pt}%
  \color{AMPrimary}\rule{\headwidth}{1.2pt}%
}
\renewcommand{\headrulewidth}{1.2pt}

% Pie: ID de reporte + paginación
\fancyfoot[L]{%
  \scriptsize\color{AMSlate}%
  ID:\ \texttt{<< report_id >>} \textbar\ v<< version >>%
}
\fancyfoot[C]{%
  \scriptsize\color{AMSlate}%
  Generado por \textbf{probe\textcolor{redtex}{.tex}} --- Una herramienta de diagnóstico de INVARIANT.%
}
\fancyfoot[R]{%
  \scriptsize\color{AMSlate}%
  Pág.\ \textbf{\thepage}\ de\ \textbf{\pageref{LastPage}}%
}
% Línea gris sobre pie
\renewcommand{\footrule}{%
  \color{AMBorder}\rule{\headwidth}{0.5pt}\vskip 2pt%
}
\renewcommand{\footrulewidth}{0.5pt}


% ════════════════════════════════════════════════════════════════════════════
%  INICIO DEL DOCUMENTO
% ════════════════════════════════════════════════════════════════════════════
\begin{document}


% ────────────────────────────────────────────────────────────────────────────
%  PORTADA
% ────────────────────────────────────────────────────────────────────────────
\thispagestyle{empty}

\noindent
{\color{AMSlate}\scriptsize\bfseries\MakeUppercase{Informe Técnico Confidencial} \hfill ID: \texttt{<< report_id >>}}\\[1ex]
{\color{AMPrimary}\rule{\linewidth}{1.5pt}}\\[4ex]

\noindent
{\color{black}\fontsize{42}{48}\selectfont\bfseries Diagnóstico}\\[1ex]
{\color{black}\fontsize{42}{48}\selectfont\bfseries Forense de Hardware} {\color{StatusCrit}\rule{14pt}{14pt}}\\[3ex]
{\color{AMSlate}\Large Hardware Audit Report}

\vspace{2.5cm}

% --- Bloques de Especificaciones ---
\noindent
\begin{minipage}[t]{0.5\textwidth}
  {\scriptsize\bfseries\color{AMSlate}HARDWARE OBJETIVO}\\[1ex]
  {\Large\bfseries\color{AMPrimary}<< device_brand >>\ << device_model >>}\\[1ex]
  {\small\color{AMSlate}S/N:\ \texttt{<< serial_number >>}}
\end{minipage}%
\begin{minipage}[t]{0.5\textwidth}
  {\scriptsize\bfseries\color{AMSlate}ENTORNO DE EJECUCIÓN}\\[1ex]
  {\normalsize\bfseries\color{AMPrimary}probe\textcolor{redtex}{.tex} v<< version >>}\\[1ex]
  {\small\color{AMSlate}Kernel Linux:\ << kernel_version >>}
\end{minipage}

\vspace{1.5cm}

\noindent
\begin{minipage}[t]{0.5\textwidth}
  {\scriptsize\bfseries\color{AMSlate}CLIENTE / PROPIETARIO}\\[1ex]
  {\normalsize\bfseries\color{AMPrimary}<< cliente_nombre >>}\\[1ex]
  {\small\color{AMSlate}<< cliente_email >>}\\[0.5ex]
  {\small\color{AMSlate}<< cliente_telefono >>}
\end{minipage}%
\begin{minipage}[t]{0.5\textwidth}
  {\scriptsize\bfseries\color{AMSlate}DIAGNÓSTICO REALIZADO POR}\\[1ex]
  {\normalsize\bfseries\color{AMPrimary}<< taller_nombre >>}\\[1ex]
  {\small\color{AMSlate}Técnico:\ << tecnico_nombre >>}\\[0.5ex]
  {\small\color{AMSlate}Fecha:\ << fecha_reporte >>}
\end{minipage}

\vspace*{\fill}

% --- Veredicto Global ---
\noindent
\begin{tikzpicture}
  \node[
    fill=AMCanvas,
    inner sep=18pt,
    text width=\textwidth - 36pt,
    draw=AMPrimary, line width=1.5pt
  ]{%
    \begin{minipage}[c]{0.65\textwidth}
      \small\color{AMText}
      \textbf{\color{AMPrimary}ANÁLISIS ESTRUCTURAL Y TERMODINÁMICO}\\[8pt]
      \textbf{Estado:}\ << estado_global_badge >>\\[6pt]
      \scriptsize\color{AMSlate}<< resumen_ejecutivo >>
    \end{minipage}%
    \hfill
    \begin{minipage}[c]{0.30\textwidth}
      \raggedleft
      {\scriptsize\bfseries\color{AMSlate}ÍNDICE DE ANOMALÍA ($\Delta A$)}\\[4pt]
      {\fontsize{32}{36}\selectfont\bfseries\color{AMPrimary}<< indice_anomalia >>}
    \end{minipage}
  };
\end{tikzpicture}

\vspace{1cm}

% --- Pie de portada ---
\begin{center}
  \large\color{AMSlate}
  Generado por \textbf{probe\textcolor{redtex}{.tex}} --- Una herramienta de diagnóstico de INVARIANT.
\end{center}

\newpage


% ────────────────────────────────────────────────────────────────────────────
%  TABLA DE CONTENIDOS
% ────────────────────────────────────────────────────────────────────────────
{
  \hypersetup{linkcolor=AMSecondary}
  \tableofcontents
}
\newpage


% ════════════════════════════════════════════════════════════════════════════
%  SECCIÓN 1 --- PROCESAMIENTO CENTRAL (CPU)
% ════════════════════════════════════════════════════════════════════════════
\section{Procesamiento Central (CPU)}

\begin{componentinfo}
  \small
  \textbf{Modelo:}\ \mv{<< cpu_model >>}\hfill
  \textbf{Núcleos/Hilos:}\ \mv{<< cpu_cores >>/<< cpu_threads >>}\hfill
  \textbf{TDP:}\ \mv{<< cpu_tdp >>\ W}\hfill
  \textbf{Socket:}\ << cpu_socket >>
\end{componentinfo}

\vspace{0.5cm}


% ────────────────────────────────────────────────────────────────────────────
%  1.1  Estabilidad Térmica Relativa
% ────────────────────────────────────────────────────────────────────────────
\subsection{Estabilidad Térmica Relativa ($\Delta T_{\mathrm{recovery}}$)}

\noindent % <--- BLOQUEA CUALQUIER SANGRÍA AUTOMÁTICA
\begin{minipage}[t]{0.58\textwidth} % <--- ACHICAMOS UN 2% PARA DAR AIRE AL MEDIO
  \vspace{0pt}
  \begin{tikzpicture}[baseline=(current bounding box.north)]
  \begin{axis}[
    probe.tex,
    width  = \linewidth,
    height = 6.2cm,
    xlabel = {Tiempo (s)},
    ylabel = {Temperatura ($^\circ$C)},
    title  = {Curva de Recuperación Térmica --- CPU},
    xmin=0, xmax=120,
    ymin=30, ymax=105,
    xtick distance=20,
    ytick distance=10,
    extra y ticks      = {<< cpu_t_idle_plus5 >>},
    extra y tick labels = {\empty}, % <--- SILENCIAMOS EL TEXTO EXTERNO
    extra y tick style  = {
      grid       = major,
      grid style = {dashed, thick, color=StatusWarnMid}
    },
    clip = false,
  ]
    % -- Curva de temperatura medida --
    \addplot[
      color=AMPrimary, very thick, smooth, mark=none
    ] coordinates { << datos_cpu_temp >> };
    \addlegendentry{$T_{\text{cpu}}$ (medida)}

    % -- Línea de T_idle de referencia --
    \addplot[
      color=StatusOKMid, thick, dashed, mark=none
    ] coordinates { (0,<< cpu_t_idle >>) (120,<< cpu_t_idle >>) };
    \addlegendentry{$T_{\text{idle}}$ = << cpu_t_idle >>$^\circ$C}

    % -- Banda de peligro TjMax (zona roja translúcida) --
    \addplot[
      fill=StatusCrit, fill opacity=0.10, draw=none
    ] coordinates { (0,<< cpu_tjmax >>) (120,<< cpu_tjmax >>) (120,110) (0,110) }
    \closedcycle;

    % -- Nodo TjMax --
    \node[
      anchor=west, font=\tiny\bfseries, color=StatusCrit
    ] at (axis cs:2, << cpu_tjmax >>+2) {TjMax (<< cpu_tjmax >>$^\circ$C)};

    % -- NUEVO NODO: T_idle + 5 (Flotando adentro del gráfico) --
    \node[
      anchor=south west, font=\tiny\bfseries, color=StatusWarn
    ] at (axis cs:2, << cpu_t_idle_plus5 >>) {$T_{\text{idle}}+5^\circ$};

  \end{axis}
  \end{tikzpicture}
\end{minipage}%
\hfill
\begin{minipage}[t]{0.39\textwidth} % <--- DAMOS UN 2% MÁS DE ANCHO A LA TABLA
  \vspace{0pt}
  \small
  \begin{tabularx}{\linewidth}{@{}X r@{}}
    \toprule
    \thcell{Parámetro} & \thcell{Valor} \\
    \midrule
    $T_{\max}$ bajo carga        & \mv{<< cpu_t_max >>$^\circ$C} \\
    \rowalt
    $T_{\text{idle}}$ medida     & \mv{<< cpu_t_idle >>$^\circ$C} \\
    $\Delta T_{\text{recovery}}$ & \mv{<< cpu_delta_t >>$^\circ$C} \\
    \rowalt
    Tiempo de recuperación       & \mv{<< cpu_recovery_time >>\ s} \\
    Límite TjMax                 & \mv{<< cpu_tjmax >>$^\circ$C} \\
    \rowalt
    Pasta térmica (estimada)     & << cpu_tim_status >> \\
    \bottomrule
  \end{tabularx}

  \vspace{0.5cm}

  [% if cpu_thermal_status == "ok" %]
  \begin{amboxok}{Estabilidad Térmica}
    \scriptsize Recuperación dentro del margen.\\
    $\Delta T =\ $\mv{<< cpu_delta_t >>$^\circ$C} (umbral: 15$^\circ$C)
  \end{amboxok}
  [% elif cpu_thermal_status == "warn" %]
  \begin{amboxwarn}{Disipación Marginal}
    \scriptsize Tiempo de recuperación elevado.\\
    Verificar pasta térmica y ventilador.
  \end{amboxwarn}
  [% else %]
  \begin{amboxcrit}{Fallo Térmico Detectado}
    \scriptsize $\Delta T$ supera umbral seguro.\\
    Intervención inmediata requerida.
  \end{amboxcrit}
  [% endif %]
\end{minipage}


% ────────────────────────────────────────────────────────────────────────────
%  1.2  Frecuencia de Throttling y Estabilidad de P-States
% ────────────────────────────────────────────────────────────────────────────
\subsection{Frecuencia de Throttling y Estabilidad de P-States}

\begin{tabularx}{\linewidth}{@{} X r r r @{}}
  \toprule
  \thcell{Modo de Operación} &
  \thcell{Reloj Base (MHz)} &
  \thcell{[% if cpu_pstate_measured %]Reloj Sostenido[% else %]Reloj Sostenido (est.)[% endif %] (MHz)} &
  \thcell{Desviación} \\
  \midrule
  P0 --- Máximo rendimiento  & \mv{<< cpu_base_p0 >>} & << cpu_sust_p0 >> & << cpu_dev_p0 >>\% \\
  \rowalt
  P1 --- Balanceado          & \mv{<< cpu_base_p1 >>} & << cpu_sust_p1 >> & << cpu_dev_p1 >>\% \\
  P2 --- Ahorro de energía   & \mv{<< cpu_base_p2 >>} & << cpu_sust_p2 >> & << cpu_dev_p2 >>\% \\
  \rowalt
  Throttling térmico (TJ90)  & ---                    & << cpu_throttle_freq >> & --- \\
  \bottomrule
\end{tabularx}

\vspace{0.3cm}

\begin{minipage}[t]{0.48\linewidth}
  \vspace{0pt} % <--- ANCLA INVISIBLE
  \small
  \textbf{Eventos de throttling:}\ \mv{<< cpu_throttle_events >>}\\
  \textbf{Duración acumulada:}\ \mv{<< cpu_throttle_dur >>\ ms}\\
  \textbf{Causa principal:}\ << cpu_throttle_cause >>
\end{minipage}
\hfill
\begin{minipage}[t]{0.48\linewidth}
  \vspace{0pt} % <--- ANCLA INVISIBLE
  [% if cpu_throttle_events == 0 %]
  \begin{amboxok}{Sin Throttling}
    \scriptsize Frecuencia plena sostenida durante toda la prueba de carga.
  \end{amboxok}
  [% else %]
  \begin{amboxwarn}{Throttling Detectado}
    \scriptsize \mv{<< cpu_throttle_events >>} evento(s). Causa: << cpu_throttle_cause >>
  \end{amboxwarn}
  [% endif %]
\end{minipage}

\vspace{0.6cm}


% ────────────────────────────────────────────────────────────────────────────
%  1.3  Machine Check Exceptions (MCE)
% ────────────────────────────────────────────────────────────────────────────
\subsection{Machine Check Exceptions (MCE)}

[% if mce_count == 0 %]
\begin{amboxok}{Sin Errores MCE Detectados}
  \small
  No se registraron Machine Check Exceptions en el log del kernel.\\[3pt]
  \textbf{Bancos MCE escaneados:}\ << mce_banks_scanned >>\quad
  \textbf{Última verificación:}\ << mce_last_check >>
\end{amboxok}
[% else %]
\begin{amboxcrit}{Errores MCE Detectados --- Fallo de Hardware Probable}
  \small
  Se detectaron \mv{<< mce_count >>} Machine Check Exception(s).\\[6pt]
  \begin{tabularx}{\linewidth}{@{} c X X l @{}}
    \toprule
    \rowcolor{StatusCritBg}
    \textbf{Banco} & \textbf{MCi\_STATUS} & \textbf{Tipo de Error} & \textbf{Timestamp} \\
    \midrule
    << mce_tabla_filas >>
    \bottomrule
  \end{tabularx}\\[4pt]
  \textbf{Recomendación:}\ << mce_recomendacion >>
\end{amboxcrit}
[% endif %]

\newpage


% ════════════════════════════════════════════════════════════════════════════
%  SECCIÓN 2 --- PROCESAMIENTO GRÁFICO (GPU)
% ════════════════════════════════════════════════════════════════════════════
\section{Procesamiento Gráfico (GPU)}

\begin{componentinfo}
  \small
  \textbf{GPU:}\ \mv{<< gpu_model >>}\hfill
  \textbf{VRAM:}\ \mv{<< gpu_vram_total >>\ GB}\ (<< gpu_vram_type >>)\hfill
  \textbf{Driver:}\ << gpu_driver_version >>\hfill
  \textbf{Bus:}\ PCIe\ << gpu_pcie_gen >>\ x<< gpu_pcie_width >>
\end{componentinfo}

\vspace{0.5cm}


% ────────────────────────────────────────────────────────────────────────────
%  2.1  Delta de Temperatura del Hotspot
% ────────────────────────────────────────────────────────────────────────────
\subsection{Delta de Temperatura del Hotspot ($\Delta T_{\mathrm{hotspot}}$)}

\begin{center}
\begin{tikzpicture}
\begin{axis}[
  probe.tex,
  xbar,
  width  = 0.88\linewidth,
  height = 4.8cm,
  xlabel = {Temperatura ($^\circ$C)},
  symbolic y coords = {Hotspot,Edge},
  ytick  = data,
  yticklabel style = {font=\small\bfseries, color=AMSlate},
  xmin=0, xmax=125,
  xtick distance=20,
  title = {Temperatura GPU: $T_{\text{edge}}$ vs $T_{\text{hotspot}}$},
  nodes near coords,
  nodes near coords style = {
    font=\scriptsize\bfseries, color=AMText,
    xshift=4pt
  },
  nodes near coords align = {horizontal},
  extra x ticks      = {<< gpu_temp_limit >>},
  extra x tick labels = {Límite \mv{<< gpu_temp_limit >>$^\circ$C}},
  extra x tick style  = {
    grid       = major,
    grid style = {dashed, very thick, color=StatusCrit, opacity=0.7},
    tick label style = {
      color=StatusCrit, font=\scriptsize\bfseries,
      rotate=90, anchor=south east, yshift=4pt
    }
  },
  legend pos = south east,
]
  % -- Barra T_edge --
  \addplot[fill=AMSlate, draw=AMSecondary, fill opacity=0.85]
    coordinates { (<< gpu_t_edge >>,Edge) };
  \addlegendentry{$T_{\text{edge}}$ = << gpu_t_edge >>$^\circ$C}

  % -- Barra T_hotspot --
  \addplot[fill=AMSecondary, draw=AMPrimary, fill opacity=0.85]
    coordinates { (<< gpu_t_hotspot >>,Hotspot) };
  \addlegendentry{$T_{\text{hotspot}}$ = << gpu_t_hotspot >>$^\circ$C}

\end{axis}
\end{tikzpicture}
\end{center}

\noindent\small
$\Delta T_{\text{hotspot}} = T_{\text{hotspot}} - T_{\text{edge}} =\ $\mv{<< gpu_delta_t_hotspot >>$^\circ$C}
\hspace{1cm}
[% if gpu_hotspot_status == "ok" %]
\badgeok\ Delta dentro del rango normal ($<$20$^\circ$C)
[% elif gpu_hotspot_status == "warn" %]
\badgewarn\ Delta elevado --- verificar pasta térmica de GPU.
[% else %]
\badgefail\ Delta crítico --- posible degradación de TIM o problema en die.
[% endif %]

\vspace{0.5cm}


% ────────────────────────────────────────────────────────────────────────────
%  2.2  Integridad de VRAM
% ────────────────────────────────────────────────────────────────────────────
\subsection{Integridad de VRAM}

[% if gpu_vram_tested %]
\begin{tabularx}{\linewidth}{@{} X c c c c @{}}
  \toprule
  \thcell{Prueba de Memoria} & \thcell{Tamaño} &
  \thcell{Patrón} & \thcell{Errores} & \thcell{Estado} \\
  \midrule
  Escritura/Lectura secuencial & \mv{<< gpu_vram_total >>\ GB} & \texttt{0xDEADBEEF} &
    \mv{<< gpu_vram_seq_errors >>} &
    [% if gpu_vram_seq_errors == 0 %]\badgeok[% else %]\badgefail[% endif %] \\
  \rowalt
  Patrón de bits aleatorio     & \mv{<< gpu_vram_total >>\ GB} & Random &
    \mv{<< gpu_vram_rand_errors >>} &
    [% if gpu_vram_rand_errors == 0 %]\badgeok[% else %]\badgefail[% endif %] \\
  Stress combinado             & \mv{<< gpu_vram_stress_gb >>\ GB} & Mixed &
    \mv{<< gpu_vram_stress_errors >>} &
    [% if gpu_vram_stress_errors == 0 %]\badgeok[% else %]\badgefail[% endif %] \\
  \rowalt
  Conteo ECC corregibles       & --- & --- & \mv{<< gpu_ecc_correctable >>} &
    [% if gpu_ecc_correctable == 0 %]\badgeok[% else %]\badgewarn[% endif %] \\
  \bottomrule
\end{tabularx}
[% else %]
\begin{amboxinfo}{Test de Integridad VRAM No Ejecutado}
  \small
  [% if gpu_vram_total == 0 %]%
    VRAM no detectada o GPU integrada sin pool dedicado.
  [% else %]%
    El análisis se ejecutó bajo entorno gráfico activo (Wayland/X11 detectado).
    Un test destructivo de escritura sobre VRAM corrompería el framebuffer
    del compositor activo.\\[4pt]
    Para ejecutar el test: reiniciar en TTY pura (\texttt{systemctl isolate multi-user.target})
    y relanzar \texttt{probe.tex}.\\[4pt]
    Herramientas requeridas según GPU:\\
    \hspace{1em}AMD\,:\enspace\texttt{yay -S gpu\_memtest}\\
    \hspace{1em}NVIDIA:\enspace\texttt{sudo pacman -S cuda-tools}\\
    \hspace{1em}Intel\,:\enspace\texttt{sudo pacman -S memtester}
  [% endif %]
\end{amboxinfo}
[% endif %]

\vspace{0.5cm}


% ────────────────────────────────────────────────────────────────────────────
%  2.3  Estabilidad del Enlace PCIe y Errores AER
% ────────────────────────────────────────────────────────────────────────────
\subsection{Estabilidad del Enlace PCIe y Errores AER}

\begin{minipage}[t]{0.56\linewidth}
  \begin{tabularx}{\linewidth}{@{} X c c @{}}
    \toprule
    \thcell{Parámetro} & \thcell{Soportado} & \thcell{Activo} \\
    \midrule
    Generación PCIe        & Gen<< gpu_pcie_gen_max >>   & \mv{Gen<< gpu_pcie_gen_active >>} \\
    \rowalt
    Ancho de banda (lanes) & x<< gpu_pcie_lanes_max >>   & \mv{x<< gpu_pcie_lanes_active >>} \\
    Ancho de banda (GB/s)  & << gpu_pcie_bw_max >>       & \mv{<< gpu_pcie_bw_active >>} \\
    \rowalt
    Errores AER corregibles & ---                         & \mv{<< gpu_aer_correctable >>} \\
    Errores AER fatales     & ---                         & \mv{<< gpu_aer_fatal >>} \\
    \bottomrule
  \end{tabularx}
\end{minipage}
\hfill
\begin{minipage}[t]{0.41\linewidth}
  [% if gpu_pcie_lanes_active == gpu_pcie_lanes_max and gpu_aer_fatal == 0 %]
  \begin{amboxok}{PCIe Estable}
    \scriptsize Enlace en capacidad máxima.\\
    Sin errores AER fatales.
  \end{amboxok}
  [% elif gpu_aer_fatal > 0 %]
  \begin{amboxcrit}{Errores AER Fatales}
    \scriptsize \mv{<< gpu_aer_fatal >>} error(es) AER fatal(es).\\
    Posible fallo en slot o conector.
  \end{amboxcrit}
  [% else %]
  \begin{amboxwarn}{Enlace PCIe Degradado}
    \scriptsize Lanes activos: x<< gpu_pcie_lanes_active >>\\
    (esperado: x<< gpu_pcie_lanes_max >>). Revisar slot.
  \end{amboxwarn}
  [% endif %]
\end{minipage}

\newpage


% ════════════════════════════════════════════════════════════════════════════
%  SECCIÓN 3 — ALMACENAMIENTO NO VOLÁTIL
%  Reemplaza completa la Sección 3 original.
%
%  ARQUITECTURA DEL LOOP
%  ─────────────────────
%  storage_list  es una list[dict] inyectada por to_jinja_context().
%  Cada `drive` es el dict plano de un StorageData individual.
%  Acceso: drive.nvme_model, drive.is_hdd, etc.
%
%  Divisor táctico al final de cada iteración excepto la última:
%  if not loop.last: (divisor) endif
% ════════════════════════════════════════════════════════════════════════════
\section{Almacenamiento No Volátil}

[% for drive in storage_list %]

\begin{componentinfo}
  \small
  \textbf{UNIDAD << loop.index >>/<< loop.length >>:}\ \texttt{<< drive.nvme_device >>}\hfill
  \textbf{Modelo:}\ \mv{<< drive.nvme_model >>}\hfill
  \textbf{Tipo:}\ [% if drive.is_hdd %]\mv{HDD Mecánico}[% else %]\mv{SSD/NVMe}[% endif %]\\[4pt]
  \textbf{Capacidad:}\ \mv{<< drive.nvme_capacity >>\ GB}\hfill
  \textbf{FW:}\ << drive.nvme_firmware >>\hfill
  \textbf{Horas de uso:}\ \mv{<< drive.nvme_hours >>\ h}
\end{componentinfo}

\vspace{0.5cm}

[% if drive.is_hdd %]
% ────────────────────────────────────────────────────────────────────────────
%  RAMA HDD: Diagnóstico de Cinemática Mecánica
% ────────────────────────────────────────────────────────────────────────────

\subsection*{Cinemática Mecánica del Cabezal y Superficie del Plato}

\begin{minipage}[t]{0.56\linewidth}
  \vspace{0pt}
  \begin{tabularx}{\linewidth}{@{} X r c @{}}
    \toprule
    \thcell{Parámetro Mecánico} & \thcell{Valor} & \thcell{Estado} \\
    \midrule
    Spin-Up Time (ID 3) & [% if drive.hdd_spin_up_time is not none %]\mv{<< drive.hdd_spin_up_time >>\ ms}[% else %]\textit{N/D}[% endif %] & --- \\
    \rowalt
    Seek Latency media & [% if drive.hdd_seek_latency_ms is not none %]\mv{<< drive.hdd_seek_latency_ms >>\ ms}[% else %]\textit{N/D}[% endif %] & [% if drive.hdd_seek_latency_ms is not none %][% if drive.hdd_seek_latency_ms > 25.0 %]\badgewarn[% else %]\badgeok[% endif %][% else %]\badgeinfo[% endif %] \\
    Sectores Reasignados & [% if drive.hdd_reallocated_sectors is not none %]\mv{<< drive.hdd_reallocated_sectors >>}[% else %]\textit{N/D}[% endif %] & [% if drive.hdd_reallocated_sectors is not none %][% if drive.hdd_reallocated_sectors > 0 %]\badgefail[% else %]\badgeok[% endif %][% else %]\badgeinfo[% endif %] \\
    \rowalt
    Command Timeouts & [% if drive.hdd_command_timeouts is not none %]\mv{<< drive.hdd_command_timeouts >>}[% else %]\textit{N/D}[% endif %] & [% if drive.hdd_command_timeouts is not none %][% if drive.hdd_command_timeouts > 0 %]\badgefail[% else %]\badgeok[% endif %][% else %]\badgeinfo[% endif %] \\
    \bottomrule
  \end{tabularx}
\end{minipage}%
\hfill
\begin{minipage}[t]{0.41\linewidth}
  \vspace{0pt}
  % --- AQUÍ ESTABA EL ERROR: Asegúrate de tener UN SOLO 'if' inicial ---
  [% if (drive.hdd_reallocated_sectors is not none and drive.hdd_reallocated_sectors > 0) or (drive.hdd_command_timeouts is not none and drive.hdd_command_timeouts > 0) %]
  \begin{amboxcrit}{Daño Físico / Fallo de Actuador}
    \scriptsize
    [% if drive.hdd_reallocated_sectors is not none and drive.hdd_reallocated_sectors > 0 %]
    \textbf{<< drive.hdd_reallocated_sectors >>} sector(es) reasignado(s). Daño físico confirmado.\\[3pt]
    [% endif %]
    [% if drive.hdd_command_timeouts is not none and drive.hdd_command_timeouts > 0 %]
    \textbf{<< drive.hdd_command_timeouts >>} timeout(s) (ID 188). Inestabilidad de actuador.\\[3pt]
    [% endif %]
    \textbf{Backup inmediato mandatorio.}
  \end{amboxcrit}
  [% elif drive.hdd_seek_latency_ms is not none and drive.hdd_seek_latency_ms > 25.0 %]
  \begin{amboxwarn}{Degradación Cinemática}
    \scriptsize Seek latency: \mv{<< drive.hdd_seek_latency_ms >>\ ms}. El cabezal pierde agilidad mecánica.
  \end{amboxwarn}
  [% else %]
  \begin{amboxok}{Cinemática Normal}
    \scriptsize [% if drive.hdd_seek_latency_ms is not none %]Seek latency: \mv{<< drive.hdd_seek_latency_ms >>\ ms}.[% endif %] Parámetros mecánicos en rango óptimo.
  \end{amboxok}
  [% endif %]
\end{minipage}

\vspace{0.5cm}

\subsection*{Métricas SMART de Superficie y Desgaste HDD}

\begin{tabularx}{\linewidth}{@{} X r @{}}
  \toprule
  \thcell{Atributo SMART} & \thcell{Valor} \\
  \midrule
  Horas de encendido (Power-On Hours)  & \mv{<< drive.nvme_hours >>\ h} \\
  \rowalt
  Spin-Up Time (ID 3, raw)             &
    [% if drive.hdd_spin_up_time is not none %]\mv{<< drive.hdd_spin_up_time >>\ ms}[% else %]\textcolor{AMSlate}{\textit{N/D}}[% endif %] \\
  Sectores Reasignados (ID 5)          &
    [% if drive.hdd_reallocated_sectors is not none %]\mv{<< drive.hdd_reallocated_sectors >>}[% else %]\textcolor{AMSlate}{\textit{N/D}}[% endif %] \\
  \rowalt
  Command Timeouts (ID 188)            &
    [% if drive.hdd_command_timeouts is not none %]\mv{<< drive.hdd_command_timeouts >>}[% else %]\textcolor{AMSlate}{\textit{N/D}}[% endif %] \\
  Seek Latency media (fio, 10 s)       &
    [% if drive.hdd_seek_latency_ms is not none %]\mv{<< drive.hdd_seek_latency_ms >>\ ms}[% else %]\textcolor{AMSlate}{\textit{N/D --- instalar fio}}[% endif %] \\
  \bottomrule
\end{tabularx}

[% else %]
% ────────────────────────────────────────────────────────────────────────────
%  RAMA SSD/NVMe
% ────────────────────────────────────────────────────────────────────────────

\begin{componentinfo}
  \small
  \textbf{TBW restante:}\ \mv{<< drive.nvme_tbw_remaining >>\ TB}\ de\ << drive.nvme_tbw_rated >>\ TB
\end{componentinfo}

\vspace{0.5cm}

\subsection*{Distribución de Latencia de Cola I/O (Histograma Logarítmico)}

\begin{center}
\begin{tikzpicture}
\begin{axis}[
  probe.tex,
  width      = 0.94\linewidth,
  height     = 6.5cm,
  ybar,
  bar width  = 14pt,
  xlabel     = {Latencia ($\mu$s)},
  ylabel     = {Operaciones (escala log)},
  ymode      = log,
  ymin       = 1,
  log origin = infty,
  title      = {Histograma de Latencia I/O --- Distribución por Bucket},
  xtick      = data,
  xticklabel style = {font=\scriptsize, rotate=40, anchor=east},
  symbolic x coords = {b0,b1,b2,b3,b4,b5,b6,b7,b8,b9},
  xticklabels = {
    ${<}1$,$1{-}2$,$2{-}4$,$4{-}8$,$8{-}16$,
    $16{-}32$,$32{-}64$,$64{-}128$,$128{-}256$,${>}256$
  },
  nodes near coords,
  nodes near coords style = {font=\tiny\bfseries, color=AMText, rotate=80, anchor=west},
  enlarge x limits = 0.07,
  legend pos = north east,
]
  \addplot[fill=AMSlate, draw=AMPrimary, fill opacity=0.82] coordinates {
    (b0,[% if drive.lat_b0 == 0 %]1[% else %]<< drive.lat_b0 >>[% endif %])
    (b1,[% if drive.lat_b1 == 0 %]1[% else %]<< drive.lat_b1 >>[% endif %])
    (b2,[% if drive.lat_b2 == 0 %]1[% else %]<< drive.lat_b2 >>[% endif %])
    (b3,[% if drive.lat_b3 == 0 %]1[% else %]<< drive.lat_b3 >>[% endif %])
    (b4,[% if drive.lat_b4 == 0 %]1[% else %]<< drive.lat_b4 >>[% endif %])
    (b5,[% if drive.lat_b5 == 0 %]1[% else %]<< drive.lat_b5 >>[% endif %])
    (b6,[% if drive.lat_b6 == 0 %]1[% else %]<< drive.lat_b6 >>[% endif %])
    (b7,[% if drive.lat_b7 == 0 %]1[% else %]<< drive.lat_b7 >>[% endif %])
    (b8,[% if drive.lat_b8 == 0 %]1[% else %]<< drive.lat_b8 >>[% endif %])
    (b9,[% if drive.lat_b9 == 0 %]1[% else %]<< drive.lat_b9 >>[% endif %])
  };
  \addlegendentry{Frecuencia de operaciones}
  \addplot[color=StatusWarn, dashed, ultra thick, mark=none, forget plot]
    coordinates { (b7, 10) (b7, 1000000) };
  \node[color=StatusWarn, font=\tiny\bfseries, rotate=90, anchor=south]
    at (axis cs:b7, 8000) {Umbral P99};
\end{axis}
\end{tikzpicture}
\end{center}

\noindent\small
\textbf{P50:}\ \mv{<< drive.nvme_lat_p50 >>\ $\mu$s}\quad
\textbf{P95:}\ \mv{<< drive.nvme_lat_p95 >>\ $\mu$s}\quad
\textbf{P99:}\ \mv{<< drive.nvme_lat_p99 >>\ $\mu$s}\quad
\textbf{P99.9:}\ \mv{<< drive.nvme_lat_p999 >>\ $\mu$s}

\vspace{0.5cm}

\subsection*{WAF y Métricas de Desgaste NVMe (SMART Extended)}

\begin{tabularx}{\linewidth}{@{} X r r c @{}}
  \toprule
  \thcell{Métrica NVMe} & \thcell{Valor Actual} &
  \thcell{Umbral} & \thcell{Estado} \\
  \midrule
  Write Amplification Factor (WAF) &
    [% if drive.nvme_waf > 0.0 %]
    \mv{<< drive.nvme_waf >>} & $\leq 3.0$ &
      [% if drive.nvme_waf_ok %]\badgeok[% else %]\badgewarn[% endif %]
    [% else %]
    \mv{N/A (no expuesto por firmware)} & --- & \badgeinfo
    [% endif %] \\
  \rowalt
  LBA escritos (host)         & \mv{<< drive.nvme_lba_written >>\ TB}  & --- & --- \\
  NAND escritas (total)       &
    [% if drive.nvme_nand_written > 0.0 %]\mv{<< drive.nvme_nand_written >>\ TB}
    [% else %]\textcolor{AMSlate}{\textit{N/D}}[% endif %] & --- & --- \\
  \rowalt
  Bloques NAND defectuosos    & \mv{<< drive.nvme_bad_blocks >>}   & $\leq 50$ &
    [% if drive.nvme_bad_blocks_ok %]\badgeok[% else %]\badgefail[% endif %] \\
  Spare blocks restantes      & \mv{<< drive.nvme_spare_blocks >>} & $\geq 10$ &
    [% if drive.nvme_spare_ok %]\badgeok[% else %]\badgewarn[% endif %] \\
  \rowalt
  Errores ECC corregibles     & \mv{<< drive.nvme_ecc_errors >>}   & $< 100$ &
    [% if drive.nvme_ecc_ok %]\badgeok[% else %]\badgewarn[% endif %] \\
  Temperatura máx. registrada & \mv{<< drive.nvme_t_max >>$^\circ$C} & $< 70^\circ$C &
    [% if drive.nvme_temp_ok %]\badgeok[% else %]\badgewarn[% endif %] \\
  \rowalt
  Porcentaje de vida restante & \mv{<< drive.nvme_life_pct >>\%} & $\geq 20\%$ &
    [% if drive.nvme_life_ok %]\badgeok[% else %]\badgewarn[% endif %] \\
  \bottomrule
\end{tabularx}

[% endif %]
% ─── FIN DE RAMA HDD / SSD ──────────────────────────────────────────────────

% ── Divisor táctico entre unidades (omitido después de la última) ───────────
[% if not loop.last %]
\vspace{0.8cm}
\noindent\textcolor{AMBorder}{\rule{\linewidth}{1pt}}
\vspace{0.5cm}
[% endif %]

[% endfor %]
% ─── FIN DEL LOOP storage_list ──────────────────────────────────────────────

\newpage


% ════════════════════════════════════════════════════════════════════════════
%  SECCIÓN 4 --- MEMORIA PRINCIPAL (RAM)
% ════════════════════════════════════════════════════════════════════════════
\section{Memoria Principal (RAM)}

\begin{componentinfo}
  \small
  \textbf{Total:}\ \mv{<< ram_total_gb >>\ GB}\hfill
  \textbf{Tipo:}\ \mv{<< ram_type >>-<< ram_speed >>}\hfill
  \textbf{Slots:}\ \mv{<< ram_slots_used >>/<< ram_slots_total >>}\hfill
  \textbf{Dual Channel:}\ [% if ram_dual_channel %]\badgeok\ Activo[% else %]\badgefail\ Inactivo[% endif %]
\end{componentinfo}

\vspace{0.5cm}


% ────────────────────────────────────────────────────────────────────────────
%  4.1  Tasa de Corrección EDAC / Bit Flipping
% ────────────────────────────────────────────────────────────────────────────
\subsection{Tasa de Corrección EDAC / Bit Flipping}

\begin{minipage}[t]{0.54\linewidth}
  \vspace{0pt} % <--- ANCLA INVISIBLE
  \begin{tabularx}{\linewidth}{@{} X c c c @{}}
    \toprule
    \thcell{Módulo DIMM} & \thcell{Tamaño} &
    \thcell{Errores} & \thcell{Estado} \\
    \midrule
    << ram_edac_filas >>
    \bottomrule
    \multicolumn{2}{@{}l}{\scriptsize\textit{Total errores detectados:}} &
    \multicolumn{2}{c@{}}{\mv{<< ram_total_errors >>}} \\
  \end{tabularx}

  \vspace{0.4cm}
  \small
  \textbf{Velocidad efectiva:}\ \mv{<< ram_speed_effective >>\ MHz}\\[2pt]
  \textbf{Latencia CAS medida:}\
  [% if ram_cas_ns is not none %]%
    \mv{<< ram_cas_ns >>\ ns}%
    \hspace{6pt}{\scriptsize\textcolor{AMSlate}{(Intel MLC \texttt{--idle\_latency})}}%
  [% else %]%
    \textcolor{AMSlate}{\textit{N/D\,---\,requiere \texttt{mlc} (Intel MLC)}}%
  [% endif %]\\[2pt]
  \textbf{Método de prueba:}\ \texttt{<< ram_test_method >>}
\end{minipage}%
\hfill
\begin{minipage}[t]{0.43\linewidth}
  \vspace{0pt} % <--- ANCLA INVISIBLE
  % -- Gráfico de pastel (donut) --- Integridad de RAM --
  \begin{center}
  \begin{tikzpicture}
    % Ángulo de la porción íntegra (inyectado por Python)
    % start: 90° (12h) en sentido antihorario
    \pgfmathsetmacro{\intAngle}{<< ram_pie_angle >>}
    \pgfmathsetmacro{\failStart}{90 - \intAngle}

    % -- Sector íntegro (verde) --
    \fill[StatusOKMid, opacity=0.9]
      (0,0) -- (90:1.85cm)
      arc[start angle=90, delta angle=-\intAngle, radius=1.85cm]
      -- cycle;

    % -- Sector defectuoso (rojo, si existe) --
    \fill[StatusCritMid, opacity=0.9]
      (0,0) -- (\failStart:1.85cm)
      arc[start angle=\failStart, delta angle=-(360-\intAngle), radius=1.85cm]
      -- cycle;

    % -- Anillo interior (efecto donut) --
    \fill[AMCanvas] (0,0) circle (1.15cm);

    % -- Texto central --
    \node[align=center] at (0,0) {%
      \textcolor{StatusOK}{\bfseries\small << ram_integrity_pct >>\%}\\[-1pt]
      \textcolor{AMSlate}{\tiny Íntegra}%
    };

    % -- Bordes --
    \draw[AMBorder, line width=0.4pt] (0,0) circle (1.85cm);
    \draw[AMBorder, line width=0.4pt] (0,0) circle (1.15cm);

    % -- Leyenda --
    \node[right, font=\scriptsize, anchor=west] at (2.05cm,  0.45cm)
      {\textcolor{StatusOKMid}{$\blacksquare$}\ Íntegra\ (<< ram_integrity_pct >>\%)};
    \node[right, font=\scriptsize, anchor=west] at (2.05cm, -0.10cm)
      {\textcolor{StatusCritMid}{$\blacksquare$}\ Defectos\ (<< ram_fail_pct >>\%)};

    % -- Título --
    \node[above, font=\small\bfseries, color=AMPrimary]
      at (0, 2.15cm) {Integridad de RAM};
  \end{tikzpicture}
  \end{center}
  \end{minipage}

  \vspace{0.6cm}

  [% if ram_total_errors == 0 %]
  \begin{amboxok}{RAM Sin Errores}
    \scriptsize Todos los módulos DIMM pasaron las pruebas de integridad.
  \end{amboxok}
  [% else %]
  \begin{amboxcrit}{Errores de Memoria Detectados}
    \scriptsize \mv{<< ram_total_errors >>} error(es) en EDAC/memtester.\\
    Aislar y reemplazar módulo defectuoso.
  \end{amboxcrit}
  [% endif %]

\newpage


% ════════════════════════════════════════════════════════════════════════════
%  SECCIÓN 5 --- PLACA BASE Y SISTEMA DE ENERGÍA
% ════════════════════════════════════════════════════════════════════════════
\section{Placa Base y Sistema de Energía}

\begin{componentinfo}
  \small
  \textbf{Placa Base:}\ \mv{<< mobo_manufacturer >>\ << mobo_model >>}\hfill
  \textbf{Chipset:}\ << mobo_chipset >>\hfill
  \textbf{BIOS/UEFI:}\ << bios_version >>\ (<< bios_date >>)
\end{componentinfo}

\vspace{0.5cm}


% ────────────────────────────────────────────────────────────────────────────
%  5.1  Caída de Voltaje bajo Carga (V_droop) y VRM
% ────────────────────────────────────────────────────────────────────────────
\subsection{Caída de Voltaje bajo Carga ($V_{\mathrm{droop}}$) y VRM}

\begin{tikzpicture}
\begin{axis}[
  probe.tex,
  width  = 0.92\linewidth,
  height = 5.8cm,
  xlabel = {Tiempo (s)},
  ylabel = {Voltaje (V)},
  title  = {VRM: Voltaje Solicitado (VID) vs Voltaje Entregado (Medido)},
  xmin=0, xmax=60,
  ymin=0.80, ymax=1.55,
  xtick distance=10,
  ytick distance=0.05,
  legend pos=south east,
]
  % -- Zona de tolerancia ±5% (relleno verde translúcido) --
  \addplot[
    fill=StatusOKMid, fill opacity=0.07, draw=none, forget plot
  ] coordinates {
    (0,  << vrm_tol_low >>)  (60, << vrm_tol_low >>)
    (60, << vrm_tol_high >>) (0,  << vrm_tol_high >>)
  } \closedcycle;

  % -- Voltaje solicitado (VID) --
  \addplot[
    color=AMSlate, thick, dashed, mark=none
  ] coordinates { << datos_vrm_vid >> };
  \addlegendentry{$V_{\text{solicitado}}$ (VID)}

  % -- Voltaje entregado (medido por sensor) --
  \addplot[
    color=AMPrimary, thick, smooth, mark=none
  ] coordinates { << datos_vrm_medido >> };
  \addlegendentry{$V_{\text{entregado}}$ (sensor)}

  % -- Anotación del droop máximo --
  \draw[->, color=AMSlate, thin]
    (axis cs:<< vrm_droop_t >>, << vrm_droop_v >>)
    -- (axis cs:<< vrm_droop_t >>+8, << vrm_droop_v >>-0.04)
    node[right, font=\tiny, color=AMSlate] {%
      $V_{\text{droop}}$\mv{= << vrm_vdroop_max >>\ V}%
    };

\end{axis}
\end{tikzpicture}

\noindent{\scriptsize\color{AMSlate}
\textit{Nota: el VRM no expone historial temporal. Las líneas representan
el voltaje VID nominal estimado y el voltaje medido puntualmente durante
el análisis (condición de carga ligera del sistema operativo).}}

\vspace{0.3cm}
\noindent\small
\textbf{$V_{\text{droop}}$ máx.:}\ \mv{<< vrm_vdroop_max >>\ V}\quad
\textbf{Temperatura VRM:}\
[% if vrm_temp is not none %]%
  \mv{<< vrm_temp >>$^\circ$C}%
[% else %]%
  \textcolor{AMSlate}{\textit{N/D}}%
[% endif %]\quad
\textbf{Fases VRM:}\
[% if vrm_phases is not none %]%
  \mv{<< vrm_phases >>}%
[% else %]%
  \textcolor{AMSlate}{\textit{N/D}}%
[% endif %]\quad
\textbf{Estado:}\ \mv{<< vdroop_status >>}

\vspace{0.3cm}
[% if vrm_status == "ok" %]
\begin{amboxok}{VRM Estable}
  \scriptsize Caída de voltaje dentro del margen de tolerancia ($\pm$5\%). Sin riesgo de inestabilidad.
\end{amboxok}
[% elif vrm_status == "warn" %]
\begin{amboxwarn}{Caída de Voltaje Moderada}
  \scriptsize $V_{\text{droop}}$ de \mv{<< vrm_vdroop_max >>\ V} detectado. Verificar fases VRM y capacitores de desacople.
\end{amboxwarn}
[% else %]
\begin{amboxcrit}{VRM Inestable --- Caída Excesiva}
  \scriptsize Caída de voltaje fuera de especificaciones. Riesgo de crashes y daño al procesador. Reemplazar VRM.
\end{amboxcrit}
[% endif %]

\vspace{0.6cm}


% ────────────────────────────────────────────────────────────────────────────
%  5.2  Estado de Batería (condicional --- solo laptops)
% ────────────────────────────────────────────────────────────────────────────
\subsection{Estado de Batería (si aplica)}

[% if battery_present %]

\begin{minipage}[t]{0.54\linewidth}
  \begin{tabularx}{\linewidth}{@{} X r @{}}
    \toprule
    \thcell{Parámetro de Batería} & \thcell{Valor} \\
    \midrule
    Capacidad de diseño         & \mv{<< bat_design_cap >>\ mWh} \\
    \rowalt
    Capacidad actual (plena)    & \mv{<< bat_full_cap >>\ mWh} \\
    Ciclos de carga             &
      [% if bat_cycles is not none %]%
        \mv{<< bat_cycles >>}%
      [% else %]%
        \textcolor{AMSlate}{\textit{N/D}}%
      [% endif %] \\
    \rowalt
    Estado de salud (SoH)       & \mv{<< bat_soh >>\%} \\
    Voltaje nominal             & \mv{<< bat_voltage_nom >>\ V} \\
    \rowalt
    Voltaje bajo carga          & \mv{<< bat_voltage_load >>\ V} \\
    $\Delta V$ bajo carga       &
      [% if bat_voltage_drop is not none %]%
        \mv{<< bat_voltage_drop >>\ V}%
      [% else %]%
        \textcolor{AMSlate}{\textit{Requiere descarga}}%
      [% endif %] \\
    \rowalt
    Corriente de descarga       & \mv{<< bat_current_load >>\ mA} \\
    Resistencia interna est.    &
      [% if bat_resistance is not none %]%
        \mv{<< bat_resistance >>\ m$\Omega$}%
      [% else %]%
        \textcolor{AMSlate}{\textit{Requiere descarga}}%
      [% endif %] \\
    \bottomrule
  \end{tabularx}
\end{minipage}%
\hfill
\begin{minipage}[t]{0.43\linewidth}
  % -- Gauge de SoH (barra de rango TikZ) --
  \begin{center}
  \begin{tikzpicture}
    % Constantes de la barra
    \def\bw{4.6}    % ancho total (cm)
    \def\bh{0.6}    % alto (cm)

    % Fondo gris
    \fill[AMBorder!60] (0,0) rectangle (\bw, \bh);

    % Relleno de color según SoH (longitud pre-calculada por Python)
    \fill[<< bat_gauge_color >>]
      (0,0) rectangle (<< bat_gauge_fill >>, \bh);

    % Divisores de zona (0-40%=critico, 40-70%=moderado, 70-100%=bueno)
    \draw[white, line width=1.2pt] (1.84, 0) -- (1.84, \bh);  % 40%
    \draw[white, line width=1.2pt] (3.22, 0) -- (3.22, \bh);  % 70%

    % Indicador de posición actual
    \draw[AMPrimary, very thick]
      (<< bat_gauge_fill >>, -0.18) -- (<< bat_gauge_fill >>, \bh+0.18);
    % Triángulo indicador
    \fill[AMPrimary]
      (<< bat_gauge_fill >>, \bh+0.18) --
      ({<< bat_gauge_fill >>-0.12}, \bh+0.38) --
      ({<< bat_gauge_fill >>+0.12}, \bh+0.38) -- cycle;

    % Etiquetas de zona
    \node[font=\tiny, color=StatusCrit]
      at (0.92, -0.28) {Crítico};
    \node[font=\tiny, color=StatusWarn]
      at (2.53, -0.28) {Moderado};
    \node[font=\tiny, color=StatusOK]
      at (3.91, -0.28) {Bueno};

    % Borde
    \draw[AMSlate, line width=0.5pt] (0,0) rectangle (\bw, \bh);

    % Valor de SoH
    \node[above, font=\bfseries\small, color=AMPrimary]
      at (\bw/2, \bh+0.55) {Estado de Salud (SoH): \mv{<< bat_soh >>\%}};

    % Título
    \node[above, font=\scriptsize, color=AMSlate]
      at (\bw/2, \bh+1.0) {Resistencia interna est.: \mv{<< bat_resistance >>\ m$\Omega$}};
  \end{tikzpicture}
  \end{center}
  \end{minipage}

  \vspace{0.6cm}

  [% if bat_soh >= 80 %]
  \begin{amboxok}{Batería en Buen Estado}
    \scriptsize SoH: \mv{<< bat_soh >>\%}. No requiere reemplazo.
  \end{amboxok}
  [% elif bat_soh >= 60 %]
  \begin{amboxwarn}{Degradación Moderada}
    \scriptsize SoH: \mv{<< bat_soh >>\%}. Considerar reemplazo próximo.
  \end{amboxwarn}
  [% else %]
  \begin{amboxcrit}{Batería Muy Degradada}
    \scriptsize SoH: \mv{<< bat_soh >>\%}. Reemplazo urgente recomendado.
  \end{amboxcrit}
  [% endif %]

[% else %]
\begin{amboxinfo}{Batería No Presente}
  \scriptsize Este sistema no cuenta con batería interna (equipo de escritorio o batería no detectada en \texttt{/sys/class/power\_supply/}).
\end{amboxinfo}
[% endif %]

\newpage


% ════════════════════════════════════════════════════════════════════════════
%  SECCIÓN 6 --- INTEGRIDAD DE BUS Y PERIFÉRICOS
% ════════════════════════════════════════════════════════════════════════════
\section{Integridad de Bus y Periféricos}


% ────────────────────────────────────────────────────────────────────────────
%  6.1  Integridad USB
% ────────────────────────────────────────────────────────────────────────────
\subsection{Integridad de Puertos USB}

\begin{tabularx}{\linewidth}{@{} c c X c c c @{}}
  \toprule
  \thcell{Puerto} & \thcell{Versión} & \thcell{Controlador} &
  \thcell{BW (Gbps)} & \thcell{Errores} & \thcell{Estado} \\
  \midrule
  << usb_tabla_filas >>
  \bottomrule
\end{tabularx}

\vspace{0.5cm}

\begin{minipage}[t]{0.48\linewidth}
  \vspace{0pt} % <--- ANCLA INVISIBLE
  \small\textbf{\color{AMPrimary}Resumen de Estado USB:}
  \begin{itemize}[leftmargin=1.2em, itemsep=3pt, topsep=5pt]
    \item \badgeok\ Puertos estables: \mv{<< usb_ok >>}
    \item \badgewarn\ Puertos inestables: \mv{<< usb_warn >>}
    \item \badgefail\ Puertos con fallo/corto: \mv{<< usb_fail >>}
  \end{itemize}
\end{minipage}
\hfill
\begin{minipage}[t]{0.48\linewidth}
  \vspace{0pt} % <--- ANCLA INVISIBLE
  [% if usb_fail > 0 %]
  \begin{amboxcrit}{Puerto(s) USB Defectuosos}
    \scriptsize \mv{<< usb_fail >>} puerto(s) con fallos.\\
    Revisar circuito de alimentación y fusibles USB.
  \end{amboxcrit}
  [% elif usb_warn > 0 %]
  \begin{amboxwarn}{Inestabilidad USB}
    \scriptsize \mv{<< usb_warn >>} puerto(s) con errores intermitentes.\\
    Verificar voltaje VBUS y controlador.
  \end{amboxwarn}
  [% else %]
  \begin{amboxok}{Todos los Puertos USB Operativos}
    \scriptsize Sin errores detectados en ningún puerto USB.
  \end{amboxok}
  [% endif %]
\end{minipage}
\newpage


% ════════════════════════════════════════════════════════════════════════════
%  TABLA RESUMEN GLOBAL + RECOMENDACIONES
% ════════════════════════════════════════════════════════════════════════════
\vspace{0.8cm}
\sectionrule

\begin{center}
  {\color{AMPrimary}\large\bfseries Resumen Global del Diagnóstico}
\end{center}
\sectionrule

\vspace{0.2cm}

\begin{tabularx}{\linewidth}{@{} X c c l @{}}
  \toprule
  \thcell{Sub-Sistema} & \thcell{$\Delta A$ Parcial} &
  \thcell{Estado Clínico} & \thcell{Anotación} \\
  \midrule
  Procesamiento Central (CPU)   & \textbf{+<< cpu_anomalia >>}  & << cpu_estado_badge >>  & --- \\
  \rowalt
  Procesamiento Gráfico (GPU)   & \textbf{+<< gpu_anomalia >>}  & << gpu_estado_badge >>  & --- \\
  Almacenamiento Interno        & \textbf{+<< nvme_anomalia >>} & << nvme_estado_badge >> & --- \\
  \rowalt
  Memoria Principal (RAM)       & \textbf{+<< ram_anomalia >>}  & << ram_estado_badge >>  & --- \\
  Placa Base / VRM Energía      & \textbf{+<< mobo_anomalia >>} & << mobo_estado_badge >> & --- \\
  \rowalt
  Bus / Periféricos (USB)       & \textbf{+<< usb_anomalia >>}  & << usb_estado_badge >>  & --- \\
  \rowalt
  Celdas de Energía (BAT)       & \textbf{+<< bat_anomalia >>}  & << bat_estado_badge >>  & --- \\
  \midrule
  \thcell{ÍNDICE DE ANOMALÍA GLOBAL} &
  \thcell{<< indice_anomalia >>} &
  \thcell{<< estado_global_badge >>} &
  \thcell{---} \\
  \bottomrule
\end{tabularx}

\vspace{0.6cm}

\begin{tcolorbox}[
  enhanced,
  colback   = AMLight,
  colframe  = AMPrimary,
  boxrule   = 0.8pt,
  arc       = 0pt,
  title     = {\bfseries Directivas de Intervención},
  coltitle  = white,
  attach boxed title to top left={yshift=-2mm, xshift=8mm},
  boxed title style={colback=AMPrimary, arc=0pt, boxrule=0pt},
  left=10pt, right=10pt, top=12pt, bottom=8pt
]
  \small
  \begin{enumerate}[leftmargin=*, itemsep=4pt, topsep=2pt]
    << lista_recomendaciones >>
  \end{enumerate}
\end{tcolorbox}


% ════════════════════════════════════════════════════════════════════════════
%  PIE DE DOCUMENTO --- FIRMA Y SELLO TÉCNICO
% ════════════════════════════════════════════════════════════════════════════
\vspace{0.8cm}
\sectionrule

\begin{minipage}[t]{0.48\linewidth}
  \small\color{AMSlate}
  \textbf{Emitido por:}\ << taller_nombre >>\\
  \textbf{Técnico certificado:}\ << tecnico_nombre >>\\
  \textbf{Fecha de diagnóstico:}\ << fecha_reporte >>\\
  \textbf{Duración del análisis:}\ << duracion_analisis >>
\end{minipage}
\hfill
\begin{minipage}[t]{0.48\linewidth}
  \raggedright\small\color{AMSlate}
  \textbf{Plataforma:}\ probe\textcolor{redtex}{.tex} v<< version >>\\
  \textbf{Kernel Linux:}\ << kernel_version >>\\
  \textbf{ID de Reporte:}\ \texttt{<< report_id >>}\\
  \textbf{Hash de integridad (SHA-256):}\ \texttt{\tiny << report_hash >>}
\end{minipage}

\vspace{0.5cm}

\begin{center}
  \scriptsize\color{AMSlate}
  Este informe es emitido con fines diagnósticos y de referencia técnica exclusivamente.\\
  Los datos fueron recolectados de forma automatizada por probe\textcolor{redtex}{.tex}
  en un entorno Live Linux offline controlado.\\[3pt]
  \textcolor{AMBorder}{\rule{0.45\linewidth}{0.3pt}}
\end{center}

\end{document}